\section{Quotient charts and the singular link}\label{sec:slices}

We construct the charts required by Section \ref{sec:comparison}.
The construction first treats a stabiliser acting freely on the
punctured zero fibre. It applies both along
$\Sigma\setminus\Delta$ and along the edge of $C_2$.
There are two parts to the argument. A slice retains the quotient
metric and all three complex structures, in addition to the moment
equations. A change of its normal coordinates then makes these
tensors agree with the flat tensors along the fixed locus. This
agreement supplies the vanishing factors needed when differentiating
the quotient metric. The algebraic identification of a quotient
alone would not supply these estimates.

\subsection{A finite-dimensional slice with its metric}
Let $\cX$ be a flat hyperk\"ahler affine space and let $\cG$ act
preserving the structure, with moment map $\mu$. We allow either
a finite-dimensional space and compact group, or the space of
unitary connections on our K3 surface. In the latter case choose
Sobolev order $s\ge3$ for connections and $s+1$ for gauge
transformations, with the $L^2$ metric, and take
$\mu(A)=(F_A^0)^+$. Fix $x_*\in\mu^{-1}(0)$, its compact
stabiliser $H$, and an $H$-invariant complement $\mathfrak q$ of
$\mathfrak h$ in the gauge Lie algebra. The four spaces
$\mathfrak qx_*,I_1\mathfrak qx_*,I_2\mathfrak qx_*,I_3\mathfrak qx_*$
are mutually orthogonal: the off-diagonal pairings are components
of the moment map evaluated on brackets, and vanish. Their
orthogonal complement is a quaternionic representation $\EE$.
Write $\EE=\EE_0\oplus\EE_1$, where $\EE_0=\EE^H$.

\begin{proposition}\label{prop:slice}
There is a real-analytic equivariant map
\[
 \iota(y)=x_*+y+c(y),\qquad
 c(y)\in\sum_a I_a\mathfrak qx_*,\quad c(y)=O(|y|^2),
\]
with $\mu^{\mathfrak q}\iota=0$. Put
$\widehat\mu=\mu^{\mathfrak h}\iota$ and $Y=\widehat\mu^{-1}(0)$.
Then $Y/H$ is a neighbourhood of $[x_*]$ in the quotient,
preserving stabilisers. There are smooth equivariant quaternionic
Hermitian tensors $(\widehat g,\widehat I_a)$ on $\EE$, equal to
the flat tensors at zero, with these properties:
\begin{enumerate}
\item on $Y$, for every $v\in T_y\EE$ and $\xi\in\mathfrak h$,
\begin{equation}\label{eq:slice-moment}
 \langle d\widehat\mu_a(v),\xi\rangle
 =\widehat g(\widehat I_a\xi y,v);
\end{equation}
\item along $\EE_0$, $\widehat\mu=d\widehat\mu=0$, and its
quadratic term in $\EE_1$ is the flat moment map for the fibre
structure $(\widehat g_u,\widehat I_a(u))|_{\EE_1}$;
\item the tensors restricted to $\EE_0$ are hyperk\"ahler, and
their quotient tensors on the free part of $Y/H$ are the original
quotient metric and complex structures.
\end{enumerate}
If $H'\subset H$ and $\mathfrak h^{H'}\subset\mathfrak h'$, then
$\widehat\mu$ vanishes on $\EE^{H'}$.
\end{proposition}
\begin{proof}
The derivative of $\mu^{\mathfrak q}$ in
$\sum I_a\mathfrak qx_*$ is invertible. In finite dimensions the
implicit function theorem and the ordinary slice theorem give
$c$ and the local quotient. For connections, write $d=d_{x_*}$.
The K3 identities and $(F_{x_*}^0)^+=0$ give, with positive
normalisation constants,
\begin{equation}\label{eq:K3-elliptic}
 (d^+)^*(\phi_a)=c_+\sum_a I_ad\phi_a,
 \qquad d^+(d^+)^*=c'_+(d^*d)\otimes1.
\end{equation}
For the off-diagonal components one uses
$d^+d\phi=[F_{x_*}^+,\phi]=0$ in the trace-free complex.
Thus $\EE=\ker d^*\cap\ker d^+$ is finite dimensional,
the cokernel is $\mathfrak h\otimes\RR^3$, and $d^*d$ is
invertible on $\mathfrak q$. Writing $c=(d^+)^*\phi$, the equation
is
\[
 \Pi_{\mathfrak q}\{d^+(d^+)^*\phi+
 ((y+(d^+)^*\phi)\wedge(y+(d^+)^*\phi))^+\}=0.
\]
It is analytic from order $s+1$ to order $s-1$, with invertible
derivative in $\phi$. The implicit function theorem and elliptic
regularity give the claimed finite-dimensional analytic map, with
values in every Sobolev order after shrinking its domain. The
Coulomb slice theorem \cite[Proposition~4.2.9]{dk90} identifies the
quotient and its stabilisers.

At $x=\iota(y)$ let $p_y$ project orthogonally off the
quaternionic span of $\mathfrak qx$. This span is invariant under
the $I_a$ even when its four summands are not orthogonal. In gauge
theory its projection is analytic: the map
\[
 Q_y(\phi_0,\phi_a)=d_x\phi_0+\sum_aI_ad_x\phi_a
\]
on $\mathfrak q\otimes\HH$ is uniformly bounded below near
$y=0$, and
$p_y=1-Q_y(Q_y^*Q_y)^{-1}Q_y^*$. The inverse follows from the
coercive form and elliptic regularity. Its complement has dimension
$\dim\EE$. Therefore $d^\perp\iota=p_y d\iota$ is invertible
onto that complement, and we define
\[
 \widehat g(v,w)=g_{\cX}(d^\perp\iota v,d^\perp\iota w),
 \qquad \widehat I_a=(d^\perp\iota)^{-1}I_a d^\perp\iota.
\]

On $Y$ the four orbit-direction summands are orthogonal. Since
$\mu^{\mathfrak q}\iota=0$, the difference
$d\iota v-d^\perp\iota v$ is in $\mathfrak qx$ there. It is
orthogonal to $I_a\mathfrak hx$. Equivariance and the ambient
moment identity give \eqref{eq:slice-moment}.
On $\cX^{H'}$ the components of $\mu$ along $\mathfrak h'$
are constant and zero. Equivariance puts their remaining values
in the dual of $\mathfrak h^{H'}$; the stated condition on $H'$
therefore forces all components to vanish. This proves the last
assertion and gives $\widehat\mu=0$ on $\EE_0$.

At $u\in\EE_0$ the linear term in the $\mathfrak h$ component
vanishes, because $\xi\iota(u)=0$ for $\xi\in\mathfrak h$.
Its quadratic term is the ambient flat moment map on
$d\iota_u\EE_1$. Replacing $d\iota_u n$ by $d^\perp\iota_u n$
does not change it. The mixed term vanishes by
$[\mathfrak h,\mathfrak q]\subset\mathfrak q$ and orthogonality
to $I_a\mathfrak qx$; the pure orbit term is a component of
$\mu(x)$ evaluated on a double bracket. This proves (2).

For the metric, minimise first over $\mathfrak qx$ and then over
$\mathfrak hx$. The result is
$\min_{\xi\in\mathfrak h}\widehat g(v-\xi y,v-\xi y)$.
If $P$ denotes the $\widehat g$-orthogonal projection off
$\mathfrak hy$, the corresponding horizontal lift is
$d^\perp\iota(Pv)$, so
the same calculation identifies the complex structures. Along
the fixed locus, the stratum is locally the smooth hyperk\"ahler
quotient of $\cX^H$ by the centraliser of $H$ modulo its ineffective
part. This proves (3).
\end{proof}

\subsection{Normalising the fibres}
We next work abstractly with the tensors and map in Proposition
\ref{prop:slice}. Suppose $H$ is free on the connected punctured
zero fibre of the flat moment map $\mu_0$ of $\EE_1$.
The fixed coordinates are $u\in\EE_0$, the normal coordinates
$z\in\EE_1$, and $\rho=|z|$.

\begin{lemma}\label{lem:normalise}
There is an equivariant $C^{1,1}$ coordinate change $\Xi$, smooth
off $\EE_0$ and equal to the identity there, carrying
$Y_0=\EE_0\times\mu_0^{-1}(0)$ onto $Y$. The pulled-back
tensors are Lipschitz, agree exactly with the flat tensors in all
normal blocks on $\EE_0$, and induce the quotient structure of
$Y/H$.
\end{lemma}
\begin{proof}
First identify the varying quaternionic Hermitian structures on
the normal fibres with a fixed one. The intertwiner
$B_u=\tfrac14(1-\sum_aI_a\widehat I_a(u))$ is invertible near
zero and commutes with $H$. Pulling the metric through $B_u$
gives a positive operator $G_u$ commuting with the $I_a$ and $H$.
Thus $A_u=G_u^{1/2}B_u$ is an equivariant quaternionic isometry
from the varying fibre to the flat fibre. In coordinates
$(u,A_u^{-1}z)$ the map becomes
\[
 \mu_1(u,z)=\mu_0(z)+R(u,z),\qquad R=O(|z|^3).
\]
Near the punctured flat zero fibre, $d\mu_0(z)$ has an equivariant
right inverse $S_z$ of homogeneity $-1$. Solve
\[
 b=S_z\{R(u,z)-\mu_0(b)\}.
\]
Contraction gives $b=O(|z|^2)$ and
$\mu_0(z+b)=\mu_1(u,z)$. After radial rescaling, differentiation
of this equation gives $|\partial^i b|\le C_i|z|^{2-i}$.
An angular invariant cutoff extends the correction off a conic
neighbourhood of the zero fibre. The estimate
$|\mu_0(z)|\ge c|z|^2\operatorname{dist}(z/|z|,\mu_0^{-1}(0))$
near that fibre shows that every sufficiently small zero of
$\mu_1$ lies where the cutoff is one. Hence the inverse of
$(u,z)\mapsto(u,z+b)$, followed by the fibre isometry, gives
$\Xi$. The second derivative bound gives $C^{1,1}$ across
$\EE_0$. The differential there is
$\id\oplus A_u^{-1}$, proving exact agreement of the normal
blocks. Equation \eqref{eq:slice-moment} ensures that the pulled-back
complex structures preserve the horizontal tangent space on $Y_0$.
\end{proof}

\subsection{Differentiating the quotient tensors}
We spell out the estimate that turns the preceding normalisation
into an adapted chart. Write $Y_0^\times=Y_0\setminus\EE_0$.
Freeness implies that this is a smooth submanifold, and the normal
space in $\EE$ is spanned by $I_a\xi z$ for $\xi\in\mathfrak h$.
It lies in $\EE_1$. If $II$ denotes its second fundamental form
for the flat metric and $\Pi_0$ the flat projection onto
$\mathfrak hz$, homogeneity and compactness of the unit link give
\begin{equation}\label{eq:singular-derivatives}
 |II|\le C/\rho,\qquad |\partial\Pi_0|\le C/\rho,
 \qquad \rho=|z|.
\end{equation}
Both $II$ and $\partial\Pi_0$ vanish in fixed directions;
their values lie in $\EE_1$. These are the only singular
coefficients in the calculation. Here $\partial$ is the flat
ambient derivative, and all estimates are on a fixed small ball,
or a closed angular subset of a free conic region.

\begin{lemma}\label{lem:basic-derivative}
Let $\overline T$ be a covariant two-tensor or an endomorphism
on $Y_0^\times/H$. Suppose its basic lift is the restriction of
an ambient tensor $T$; for an endomorphism require that $T$
maps $TY_0^\times$ into itself. Let $N(T)$ be the supremum,
over unit normal vectors $\nu$, of
$|T(\nu,\cdot)|+|T(\cdot,\nu)|$ in the first case, and of
$|T\nu|$ in the second. Then
\begin{equation}\label{eq:basic-derivative}
 |\nabla^{\bar g_{\fl}}\overline T|
 \le C\bigl(|\partial T|+|II|N(T)\bigr).
\end{equation}
The constant depends only on the dimension. A basic lift is zero
on orbit directions and, for an endomorphism, takes values in the
flat horizontal space.
\end{lemma}
\begin{proof}
For horizontal lifts $v,w$ of vector fields on the quotient,
the horizontal part of $\nabla^{Y_0}_vw$ is the lift of the
quotient covariant derivative. Vertical components disappear on
evaluating the basic tensor. For a covariant tensor, the Gauss
formula gives
\[
 (\nabla^{Y_0}_vT)(w_1,w_2)
 =(\partial_vT)(w_1,w_2)
   +T(II(v,w_1),w_2)+T(w_1,II(v,w_2)).
\]
For an endomorphism, project the corresponding formula to
$TY_0^\times$. The additional term is the tangential projection
of $TII(v,w)$; the term $II(v,Tw)$ has zero tangential
projection. Restricting to horizontal vectors and taking norms
proves the estimate.
\end{proof}

In particular, a bounded ambient derivative is sufficient on the
quotient when the tensor's normal components are $O(\rho)$.
The curvature of the zero fibre in \eqref{eq:singular-derivatives}
then contributes a bounded term.

\begin{lemma}\label{lem:quotient-estimate}
Suppose invariant Lipschitz tensors $(\widehat g,\widehat I_a)$
on $Y_0$ arise as above, agree with the flat tensors in normal
blocks along $\EE_0$, and preserve the horizontal tangent spaces.
Their quotients differ by $O(\rho)$ from the product of their
restrictions to $\EE_0$ with the flat cone quotient, with bounded
covariant first derivative. If two such sets of tensors differ
by $O(t)$ in Lipschitz norm and share the normal blocks, then their
quotient difference, after subtracting the difference on the
smooth factor, is $O(t\rho)$ with derivative $O(t)$.
\end{lemma}
\begin{proof}
\emph{Normal components.}
Put $P_0=1-\Pi_0$, $\tau=\widehat g-g_{\fl}$ and
$\tau_a=\widehat I_a-I_a$. Invariance makes the fixed and
normal representations orthogonal along $\EE_0$ and makes the
complex structures preserve both. Normalisation and the Lipschitz
bound therefore give, for $v_1\in\EE_1$,
\begin{equation}\label{eq:normal-errors}
 \begin{split}
 |\tau(v_1,\cdot)|\le C\rho|v_1|,\quad
 |\tau_av_1|\le C\rho|v_1|,\quad
 |(\tau_av)_1|\le C\rho|v|.
 \end{split}
\end{equation}
The fixed--fixed components need not vanish: their restrictions
are the metric and complex structures of the smooth stratum.

\emph{Orbit projections.}
Let $\Pi$ be the $\widehat g$-orthogonal projection onto
$\mathfrak hz$ and $P=1-\Pi$. With
$\tau^\sharp=g_{\fl}^{-1}\tau$, set
$A=1+\Pi_0\tau^\sharp\Pi_0$ on the ambient space. This
is invertible near the fixed locus. The equation
$\Pi_0(1+\tau^\sharp)(v-\Pi v)=0$ implies
\begin{equation}\label{eq:projection-error}
 \Pi-\Pi_0=
 A^{-1}\Pi_0\tau^\sharp P_0.
\end{equation}
Its size is $O(\rho)$ by \eqref{eq:normal-errors}. To
differentiate it, consider the factors in turn. The term
$(\partial\Pi_0)\tau^\sharp P_0$ is bounded because
$\partial\Pi_0$ kills $\EE_0$ and the normal component of
$\tau^\sharp P_0$ is $O(\rho)$. The term
$\Pi_0(\partial\tau^\sharp)P_0$ is bounded by the Lipschitz
assumption. The term $\Pi_0\tau^\sharp\partial P_0$ is
bounded by $C\rho\cdot C/\rho$. Finally differentiating
$A^{-1}$ contributes at worst $C/\rho$, again multiplied
by $\Pi_0\tau^\sharp P_0=O(\rho)$. We obtain
\begin{equation}\label{eq:projection-bounds}
 |P-P_0|\le C\rho,\qquad |\partial(P-P_0)|\le C.
\end{equation}

\emph{The quotient metric.}
The basic lift of its difference from the flat quotient metric is
the restriction to $TY_0^\times$ of
\[
 T=g_{\fl}(P\cdot,P\cdot)-g_{\fl}(P_0\cdot,P_0\cdot)
       +\tau(P\cdot,P\cdot).
\]
Since $P-P_0$ takes values in $\EE_1$, $Pv$ and $v$ have
the same fixed component $v_0$. At $(u,z)$ the tensor
$\tau$ is $\tau_u+O(\rho)$, and $\tau_u$ has only a
fixed--fixed block. Hence
\[
 T(v,w)=(g_B-g_{\fl})(u)(v_0,w_0)+O(\rho)|v||w|.
\]
This identifies the product term that must be subtracted. To
bound $\partial T$, expand $P=P_0+(P-P_0)$ in the first
two terms. Derivatives of $P-P_0$ are bounded; every derivative
of $P_0$ is multiplied by an $O(\rho)$ factor. In the last
term, $\partial\tau$ is bounded, while the potentially
singular part of $\partial P$ has normal image and is absorbed
by \eqref{eq:normal-errors}. Thus $|\partial T|\le C$.
For a unit normal $\nu$ to $Y_0^\times$, $P_0\nu=\nu$
and $P\nu\in\EE_1$, so $|T(\nu,\cdot)|\le C\rho$.
Lemma \ref{lem:basic-derivative} now gives the bounded
covariant derivative on the quotient.

\emph{The complex structures.}
Their basic lifts, in the flat horizontal identification of the
quotient tangent space, are $P_0\widehat I_aP$ and
$P_0I_aP_0$. Their difference is the restriction of
\[
 J_a=P_0\tau_aP+P_0I_a(P-P_0).
\]
Here $(\partial P_0)\tau_aP$ is bounded by the last estimate
in \eqref{eq:normal-errors}; $P_0(\partial\tau_a)P$ is
bounded by Lipschitz regularity; and $P_0\tau_a\partial P$
is bounded because the singular part of $\partial P$ takes
values in $\EE_1$. The remaining terms are controlled by
\eqref{eq:projection-bounds}. For unit normal $\nu$,
$|J_a\nu|\le C\rho$. The basic-derivative lemma applies
again. The value of $J_a$ is the fixed-factor difference
$(I_{a,B}-I_a)(u)$, extended by zero on $\EE_1$, up to
$O(\rho)$. This proves the absolute estimates for the metric
and all three complex structures.

\emph{Dependence on the data.}
For the relative assertion interpolate the two metrics by
$g_s=g_*+s\tau_*$. If $P_s$ is its horizontal projection,
stationarity of the minimum over orbit directions gives
\[
 \frac{d}{ds}g_s(P_sv,P_sw)=\tau_*(P_sv,P_sw).
\]
Indeed $\partial_sP_sv$ is an orbit vector and is
$g_s$-orthogonal to $P_sw$, so the derivatives of the two
projections contribute zero to this formula. Integration in $s$
expresses the difference of the quotient metrics as the integral
of the displayed tensor. Its fixed-factor restriction is
$\tau_*|_{\EE_0}$, and its normal blocks are $O(t\rho)$.
Every differentiated term just estimated is therefore $O(t)$,
including the second-fundamental-form contribution.

For the complex structures use
\[
 \frac{d}{ds}\Pi_s=\Pi_sg_s^{-1}\tau_*P_s
\]
and integrate to obtain
$|P_{s=1}-P_{s=0}|\le Ct\rho$ and
$|\partial(P_{s=1}-P_{s=0})|\le Ct$. The difference of
the complex-structure lifts is
\[
 P_0(\widehat I_a-\widehat I_{a,*})P_{s=1}
       +P_0\widehat I_{a,*}(P_{s=1}-P_{s=0}).
\]
The first term has normal vanishing $O(t\rho)$ by the
common normalisation of the complex structures, and the second
has that bound by the projection formula. Their derivatives are
$O(t)$ by the same cancellations. After subtracting the
fixed-factor restriction, the difference is $O(t\rho)$.
This proves the relative assertion.
\end{proof}

Together, the two lemmas construct adapted charts for each complex
structure, with model the actual stratum metric times the flat
transverse cone. The quotient map is a homeomorphism for quotient
topologies, smooth on regular loci. Its metric comparison also
identifies completions. The smooth factor may be completed to a
compact K\"ahler manifold as in Section \ref{sec:comparison}.

\subsection{The link $Z_2$ and induction on depth}\label{ss:link-induction}
Apply these charts first in finite dimensions to the rank-one
stratum of $C_2$, using \eqref{eq:rank-one-slice}, and then in
gauge theory to $\Sigma\setminus\Delta$. In both cases the
normal circle action is free off zero and the local model is
$B\times C_1$.

\begin{proposition}\label{prop:models}
The link $Z_2$ is admissible. The Stokes property holds on $M$,
and the Gaffney domain is Weil invariant on $C_2$ and on
$M\setminus\Delta$, for all three complex structures.
The models $B\times C_2$ have bounded geometry at scale $\rho$
and uniform edge charts.
\end{proposition}
\begin{proof}
Stokes is invariant under quasi-isometry, since the minimal and
maximal graph norms of $d$ are equivalent. It is local for forms
supported away from the artificial boundary of a chart, and on a
compact space these local statements combine by Lipschitz
partitions of unity. The odd smooth link $Z_1$ gives Stokes on
$C_1$; the product lemma gives it on $B\times C_1$. The adapted
charts therefore give Stokes on $C_2\setminus\{0\}$ locally
and on $M\setminus\Delta$ locally. Annular cone coordinates
are quasi-isometric to interval products with $Z_2$. Applying
local Stokes on a circle product and then to forms pulled back
from $Z_2$ proves Stokes on $Z_2$.

To prove compact resolvent, take a sequence bounded in $V(Z_2)$
and pull it back to the cone, multiplying by a fixed radial cutoff
in $1<r<2$. The cone differential formulas give a bounded
sequence in $V(C_2)$ with fixed compact support. A finite cover
by adapted $B\times C_1$ charts transfers its graph norm to
product graph norms. Lemma \ref{thm:product} makes each local
piece precompact. Its cone $L^2$ norm is a fixed positive
multiple of its link norm, degree by degree. Thus
$V(Z_2)\hookrightarrow L^2(Z_2)$ is compact. Oddness of
$\dim Z_2$ now gives Stokes on $C_2$, proving admissibility.
The imported quasi-isometry \eqref{eq:imported-metric} transfers
Stokes from $B\times C_2$ to $M$ near $\Delta$; compactness
of $M$ completes the global assertion.

For Weil invariance, use the smooth-link cone theorem and the
product lemma on $B\times C_1$, then Theorem
\ref{thm:comparison} in the adapted charts. This treats
$C_2\setminus\{0\}$ and $M\setminus\Delta$. Theorem
\ref{thm:vertex} applies to the now admissible $Z_2$ and treats
the vertex of $C_2$.

Finally the distance $\rho$ on $C_2$ is homogeneous. On the
unit annulus, finitely many adapted edge charts and smooth
charts supply bounded geometry; dilations give it at all
scales. For doubling, balls of radius comparable to their
distance from the vertex have volume comparable to the
appropriate power of their radius, by homogeneity and compactness
of the link. Smaller balls reduce to the unit annulus and the
doubling product charts. Edge charts on that annulus dilate to
uniform edge charts at every radius. Taking products with
small balls in $B$, extended to compact smooth factors, proves
the assertion for $B\times C_2$.
\end{proof}
